\centerline{\bf \Huge Постулат Бертрана}

\begin{flushright}
        \scriptsize{Эпизод 21. Рождение Nerv}
\end{flushright}

Обозначим за $R(n)$ произведение всех простых чисел от $n+1$ до $2 n$ (если таких нет, считаем произведение равным 1). Обозначим за $\varphi(n)$ количество простых чисел, не больших $n$.

\problem{1}
Докажите, что:

(a) если $2 \leqslant 2 n / 3<p \leqslant n$, то $C_{2 n}^n$ не делится на $p$.

(b) если $p>\sqrt{2 n}$, то $C_{2 n}^n$ не делится на $p^2$.

(c) если $C_{2 n}^n$ делится на $p^k$, то $p^k \leqslant 2 n$.

\problem{2}
Докажите, что $C_{2 n}^n \geqslant \frac{4^n}{2 \sqrt{n}}$.

\problem{3}
Докажите, что произведение всех простых чисел от 1 до $n$ не превосходит $4^n$.

\problem{4}
Докажите, что $R(n)>\frac{4^{n / 3}}{2 \sqrt{n}(2 n)^{\varphi(\sqrt{2 n})}}$.

\problem{5}
Докажите, что:

(a) $\varphi(n) \leqslant n / 2$.

(b) $(2 n)^{\sqrt{n / 2}}<2^{n / 3}$ для $n>500$.

(c) $R(n)>1$ для $n>500$.

\problem{6} 
Докажите постулат Бертрана.